\section{微扰论}\label{sec:perturbation}

场 $B^a_{\mu}(t,x)$ 的 $n$ 点关联函数可以按微扰论计算：如第 \ref{sec:iterative} 节
所解释的，把该场按基本场的幂次展开，
然后照常在 $\SUn$ 规范理论中计算后者的关联函数。
现在我们证明，这些关联函数可以直接由 $D+1$ 维中的一组费曼规则得到。

\subsection{规范固定}\label{ssec:gaugefixing}

流方程 \eqref{eq:2.4} 在无穷小变换
\begin{equation}
  \dbrs B_{\mu}=D_{\mu}\omega
  \label{eq:3.1}
\end{equation}
下不变，只要 $\omega(t,x)\in\sun$ 的时间依赖性满足
\begin{equation}
  \partial_t\omega=\alpha_0D_{\mu}\partial_{\mu}\omega.
  \label{eq:3.2}
\end{equation}
由于 $\omega$ 的初值不受约束，这些变换在流时间 $t=0$ 处生成整个规范群，
从而把 $\SUn$ 规范理论的规范对称性拓展到所有流时间。

可以照常通过加入规范固定项
\begin{equation}
  \Sgf=-{\lambda_0\over g_0^2}\int\rmd^Dx\,
  \tr\{\partial_{\mu}A_{\mu}(x)\partial_{\nu}A_{\nu}(x)\}
  \label{eq:3.3}
\end{equation}
及相应的鬼场作用量
\begin{equation}
  \Sgh=-{2\over g_0^2}\int\rmd^Dx\,
  \tr\{\partial_{\mu}\cbar(x)D_{\mu}c(x)\}
  \label{eq:3.4}
\end{equation}
来固定该对称性，二者计入理论的总作用量，其中 $c$ 与 $\cbar$
是 Faddeev--Popov 鬼场。就基本规范场与鬼场的 $n$ 点关联函数而言，
费曼规则便是标准规则。显然，既然变换 \eqref{eq:3.1}
是无穷小规范变分，流生成的场中规范不变表达式的期望值
便不依赖规范参数 $\lambda_0$。

\subsection{规范场传播子}

流线图如何与底层理论的费曼图相结合，最好通过考虑随时间变化的规范场的
两点函数来说明。在规范耦合的领头阶，流线图
\begin{equation}
  \raisebox{-0.5cm}{\includegraphics[width=3.5cm]{images/diagram2}}
  \label{eq:3.5}
\end{equation}
是对该关联函数有贡献的唯一图形。一点顶点处规范场的收缩表明
\begin{align}
  \langle \tilde{B}^a_{\mu}(t,p)\tilde{B}^b_{\nu}(s,q)\rangle=
  (2\pi)^D\delta(p+q)\delta^{ab}g_0^2\Dtilde_{t+s}(p)_{\mu\nu}+\rmO(g_0^4),
  \label{eq:3.6}\\[2pt]
  \Dtilde_t(p)_{\mu\nu}=
  {1\over(p^2)^2}
  \bigl\{
  (\delta_{\mu\nu}p^2-p_{\mu}p_{\nu})\rme^{-tp^2}+
  \lambda_0^{-1}p_{\mu}p_{\nu}\rme^{-\alpha_0tp^2}\bigr\}.
  \label{eq:3.7}
\end{align}
这一公式包含了混合传播子
\begin{equation}
  \langle \tilde{A}^a_{\mu}(p)\tilde{B}^b_{\nu}(s,q)\rangle=
  \langle \tilde{B}^a_{\mu}(0,p)\tilde{B}^b_{\nu}(s,q)\rangle
  \label{eq:3.8}
\end{equation}
以及流时间为零处的规范场两点函数。

既然所有三个传播子都由同一个解析表达式给出，
便可用同一个图形符号
\begin{equation}
  \raisebox{-0.55cm}{\includegraphics[width=2.5cm]{images/prop2}}%
  =\;\delta^{ab}g_0^2\Dtilde_{t+s}(p)_{\mu\nu}
  \label{eq:3.9}
\end{equation}
表示它们。注意，一点顶点的收缩总是起到把终端流线转换为规范场线的作用。
例如，若从流线图
\begin{equation}
  \raisebox{-0.6cm}{\includegraphics[width=4.5cm]{images/diagram4}}
  \label{eq:3.10}
\end{equation}
而非图 \eqref{eq:3.5} 出发，则领头阶贡献可以通过把三个一点顶点处的
规范场关联函数换成树级图而得到。这样便得到一个图
\begin{equation}
  \raisebox{-0.5cm}{\includegraphics[width=3.6cm]{images/diagram5}}
  \label{eq:3.11}
\end{equation}
它含有一个普通三点顶点（实心圆点）和一个流顶点。
连接到后者的线中有两条是规范场线而非流线，
但顶点的表达式与前相同。

\subsection{$D+1$ 维中的费曼规则}

现在把流时间解释为一个附加的时空坐标。由于只考虑非负时间，
这个 $D+1$ 维空间是一个半空间，其在流时间为零处有一个 $D$ 维边界。
$\SUn$ 规范理论生活在边界上，而由梯度流生成的场延伸到额外维中。

从这一视角看，普通顶点与流顶点分别代表边界相互作用项与体相互作用项，
而场在 $D+1$ 维中的传播由规范场传播子 \eqref{eq:3.9} 与流传播子 \eqref{eq:2.16}
描述。另一方面，鬼场传播子
\begin{equation}
  \langle\ctilde^a(p)\cbartilde^b(q)\rangle=
  (2\pi)^D\delta(p+q)\delta^{ab}g_0^2\Dtilde(p)+\rmO(g_0^4),
  \qquad
  \Dtilde(p)={1\over p^2},
  \label{eq:3.12}
\end{equation}
只在流时间为零处有定义。图形上它用一条虚线（dotted line）表示：
\begin{equation}
  \raisebox{-0.55cm}{\includegraphics[width=1.8cm]{images/prop3}}%
  =\;\delta^{ab}g_0^2\Dtilde(p),
  \label{eq:3.13}
\end{equation}
箭头沿从 $\cbar$ 指向 $c$ 的方向画出。

\begin{figure}[t]
\centerline{\includegraphics[width=10.0cm]{images/diagram6}}
\caption{对随时间变化的规范场的两点与三点函数有贡献的图的示例。
所有的图都由下列元素构成：流传播子（带方向的实线）、
规范场传播子（波浪线）、鬼场传播子（带方向的点线）、
流顶点（空心圆圈）以及流时间为零处的普通顶点（实心圆点）。
外线末端的小方块表示这些线未被截断。}
\label{fig:examples}
\end{figure}

至此应当相当清楚：梯度流生成之场的 $n$ 点函数
由 $D+1$ 维中的费曼图给出，其图形元素列于图~\ref{fig:examples} 的图注中。
费曼规则的若干特点值得指出：

\begin{enumerate}[label=(\alph*),itemsep=1.0ex,leftmargin=*]
\item 每个流顶点恰有一条朝外的流线，对应于式 \eqref{eq:2.13} 中
第一个动量--指标组合。连接到这些顶点的其他线可以是规范场线或朝内的流线。

\item 流线必须从某个流顶点出发，或者成为外线，或者终止于另一个流顶点。
而规范场线既可以从流顶点出发、也可以从普通顶点出发，
并终止于这两类顶点。

\item 流顶点插入在某个从零积到无穷大的流时间处；
所有其他顶点都位于流时间为零处。传播子所依赖的流时间
就是相应线两端点处的流时间。

\item 带闭合流线圈的图一律取为零。这条规则源于如下事实：
流线图是树图，而且一点顶点处规范场的收缩决不会产生新的流线。
\end{enumerate}
